% ------------------------------------------------------------
% Page 111
% ------------------------------------------------------------

\section*{FERRUSSAC}

\subsection*{Un village exemplaire ?}

La vie communautaire y était particulièrement développée et nous disposons de plusieurs souvenirs et
documents qui l’attestent.

\textbf{Projet de temple}

Après la reconstruction du temple de Meyrueis en 1840, « les campagnes, en
grande majorité huguenotes, écrit Philippe Chambon, veulent se doter de lieux de culte. La commune
de Gatuzieres s’engage, avec l’autorisation de l’Etat dans l’édification d’un temple, toujours debout.
Vers la même époque, les habitants des Oubrets, à 11km de Meyrueis sont autorisés à aménager un
lieu de culte. Ce temple est fort original car il a une double fonction : en semaine, il abrite l’école
communale du hameau et deux dimanches par mois, ainsi qu’aux fêtes il retrouve sa fonction cultuelle.
Cette situation se maintiendra jusqu’en 1945 époque de la fermeture de l’école et de la suppression
des cultes aux Oubrets. Ce petit temple est aujourd’hui fort bien rénové en maison d’habitation ».

Les habitants de Ferrussac, eux aussi, voulurent construire leur temple, mais l’autorisation leur fut
refusée sans doute à cause de la proximité de Meyrueis (3km seulement). Peu importe, une salle fut
aménagée, permettant des veillées communes et l’ « écolette », école biblique du jeudi pour les enfants
du hameau, rassemblés autour du pasteur. Cette salle existe toujours, salle communautaire qui peut
abriter rencontres et réunions diverses, politiques ou sociales, comme du matériel agricole appartenant
à la section communale.

\textbf{La scie :}

L’équipement agricole pouvait être acheté avec les fonds de la section, notamment les coupes
de bois. C’est ainsi que fut financée une scie circulaire avec son chariot, permettant de débiter, sur
place, planches, poutres et chevrons.

\textbf{Le troupeau de Ferrussac}

Chaque matin un des habitants, à tour de rôle, rassemblait les chèvres de
toutes les maisons et ce troupeau du village, passant d’une parcelle à l’autre, n’employait qu’un seul
berger en l’occurrence souvent une seule bergère.

\textbf{L’ancien four à pain et le moulin}

étaient tombés en ruines à la suite des inondations catastrophiques
de 1900 et peut-être par manque d’entretien régulier. C’est seulement en 1917-1918 qu’ils seront
remis à neuf. La reconstruction du Moulin s’élèvera à 3.920 F et celle du Four 1.330 F
« ce qui fait une somme globale de 5250F, laquelle dépense sera couverte par la vente d’arbres,
essence pin, (sic) qui appartiennent à la section de Ferrussac et qui sont soumis au régime forestier. Le
conseil, après examen, reconnaissant l’utilité et même l’urgence qu’il y a de faire exécuter ces travaux,
le plus tôt possible, pour que les habitants ne soient point obligés de venir moudre leurs grains et faire
leur pain à Meyrueis, à l’unanimité, donne un avis favorable à la demande des habitants de Ferrussac
(délibération du Conseil Municipal de Meyrueis. 10 juin 1917).
